\subsection{Folding an arbitrary pendant stalk to a fixed arrow}

\begin{theorem}[Proved: an acyclic pendant corner and a fixed retained arrow]
\label{CP2:A-pendant-arbitrary-stalk-fixed-fold}
Let $p$ be an object of a finite category $C$ with exactly one strict
neighbor $q$. Suppose $C(p,q)$ is nonempty and $C(q,p)$ is empty.
Put $L=C(p,p)$ and $N=C(q,q)$, and assume that $BL$ is rationally
acyclic. Suppose a retained object $r\ne p,q$ has an arrow
$g:r\to q$ such that $ng=g$ for every $n\in N$.
Let $D$ be the full subcategory obtained by deleting $p$.
If $BC\simeq\CP^2$, then $BD\simeq\CP^2$.
The opposite statement holds, with a retained arrow $g:q\to r$
satisfying $gn=g$ for every $n\in N$.
There is no bound on the nonempty pendant stalk $C(p,q)$.
\end{theorem}
\begin{proof}
Write $S=C(p,q)$. Define a functor $R:C\to D$ by fixing $D$,
sending $p$ to $r$, sending every element of $L$ to $1_r$,
and sending every arrow of $S$ to $g$.
The products involving an arrow of $S$ have the form $ns\ell$,
with $n\in N$, $s\in S$, and $\ell\in L$.
Their images are all $g$, and $ng1_r=g$ by hypothesis.
The corner map $L\to\{1_r\}$ is unital and multiplicative.
These are all new composition equations, since $p$ has no other
strict neighbor. In particular $R$ is an actual functor, and
its composite with the full inclusion $i:D\hookrightarrow C$
satisfies $Ri=1_D$.

Here is the relative product calculation without imposing any
condition on the actions on $S$. The actual two-block bar model
for the cut $\{p\}\mid D$ is the homotopy pushout
\[
 BL\ \longleftarrow\ E=B(N,S,L)\ \longrightarrow\ BD.
\]
Indeed no strict arrow can leave $q$ for an object other than $p,q$:
composing it with an element of $S$ would give another strict
neighbor of $p$. Hence the bridge has just the displayed $N$ and
$L$ bar directions, and its right map factors through $BN\subset BD$.
Use the double mapping cylinder, keeping $BD$ as a subcomplex.
The relative quotient is the homotopy cofiber of $E\to BL$.
Collapsing its copy of $BL$ gives the unreduced suspension
$\Sigma E$. Since $BL$ is connected and rationally acyclic,
this collapse induces an isomorphism on rational reduced
homology and therefore on rational cohomology rings.
Every positive-degree relative rational product in
$H^*(BC,BD;\mathbb Q)$ is consequently zero.
All the comparisons are actual maps of the bar pairs, so this
is the relative product for the original inclusion of nerves.

Assume now $BC\simeq\CP^2$. The actual retraction makes
$H_j(BD;\mathbb Z)$ a direct summand of $H_j(BC;\mathbb Z)$,
and makes $\pi_1(BD)$ a retract of the trivial group.
If $H^2(BD;\mathbb Q)=0$, the degree-two generator of
$H^*(BC;\mathbb Q)$ lifts to relative cohomology and has zero
square by the preceding paragraph, a contradiction.
Thus $H_2(BD;\mathbb Z)=\mathbb Z$, and the split degree-two
map induced by $R$ is an isomorphism. Its generator has square
pulling back to the primitive degree-four generator of $BC$.
The split degree-four map is therefore also an isomorphism.
All other integral homology groups agree as well, since they
are direct summands of the corresponding zero groups of $BC$.
The spaces are simply connected CW spaces, so the resulting
integral homology equivalence is a homotopy equivalence.
This proves $BD\simeq\CP^2$. Opposition gives the dual case.
\end{proof}

\begin{corollary}[Proved: a cardinality test at nineteen arrows]
\label{CP2:A-nineteen-pendant-arbitrary-stalk-fold}
A nineteen-arrow category cannot have type $\CP^2$ if it has
a source or sink pendant $p$ whose corner has order at most four,
and another retained singleton hom-set into its neighbor $q$
in the source case, or out of $q$ in the sink case.
The pendant stalk may have any positive size.
\end{corollary}
\begin{proof}
Corner classifying spaces through order four are rationally
acyclic. The unique arrow in the retained singleton hom-set
is fixed by every endomorphism of $q$ on the required side.
The theorem gives a projective-plane model after deleting $p$.
The deletion removes its nonempty corner and its nonempty stalk,
so has at most seventeen arrows, contradicting
Theorem~\ref{CP2:A-global-lower-nineteen}.
\end{proof}
